\begingroup
\raggedbottom

\subsubsection{Candidate-Specific Effects and Full-Selector Quality}
\label{app:rebuttal_selector}

This study extends Table~\ref{tab:effect_probes} and Appendix~\ref{app:effect_probes} from effect accuracy to decision quality.
The central question is whether $\widehat E_i$ distinguishes actions within the same candidate pool and improves the deployed selector, rather than only predicting a query's average future.
The expanded evaluation uses a separate cohort and complements, rather than replaces, the main-text diagnostic.

\paragraph{Shared queries, branches, and outcome labels.}
The cohort comprises 20 held-out source episodes per task on Put Back Block (PBB), Press Button (PB), and Battery Try (BT), with ten prespecified replans per episode and the deployed $K=32$ candidate slots at each replan.
Episodes used for fitting or hyperparameter/checkpoint selection are excluded.
The query schedule is fixed before inspecting predictions: five replans come from approach/alignment stages and five from contact/completion stages of the relevant interaction, using simulator task predicates rather than model scores to identify stages.
Episodes with fewer available queries are retained with their actual counts; they are not replaced according to outcomes.
At each query, we cache the requirement, adapted proposals, rank scores, and cursor state once.
We execute each proposal for $H=32$ commands from the same restored simulator, controller, task-progress, RNG, and factual-history state, following Appendix~\ref{app:effect_probes}; no Action DiT refinement is applied in this branch diagnostic.

Applicability is a task-progress label, not a threshold on $\kappa_i$.
For PBB, a branch must advance the correct grasp/return interaction toward the episode's original pad, without placing the block on a wrong pad.
For PB, it must advance approach to the required button or register a permitted press without exceeding the remaining count.
For BT, it must advance alignment/insertion of the battery currently required by task feedback, without repeating a feedback-rejected attempt.
Before evaluation, each sampled stage is assigned one explicit predicate: reaching the task evaluator's contact/alignment region, acquiring the correct object, or completing the required placement/press/insertion.
The mapping and geometric tolerances are frozen on development episodes; a branch is positive only if that predicate is reached without an invalid transition.
No progress or an incorrect interaction gives $a_{t,i}=0$; an outcome that cannot be resolved from the simulator state is unknown, not automatically negative.
Labels use privileged state only for evaluation, never as policy or selector inputs.

\begin{table}[H]
\centering\small
\caption{\textbf{Branch-cohort accounting.}
Full endpoints, early task termination, and missing endpoints partition attempted branches.
Applicability counts are separate: an early terminal outcome can still have a known label.
Coverage and choice accuracy use the same fully labelled query cohort.}
\label{tab:rebuttal_coverage}
\begin{tabular}{@{}lrrr@{}}
\toprule
Accounting item & PBB & PB & BT \\
\midrule
Source episodes / queries & 20 / 200 & 20 / 200 & 20 / 200 \\
Attempted candidate branches & 6,400 & 6,400 & 6,400 \\
Complete $H$-step endpoints & 6,304 & 6,272 & 6,208 \\
Early task termination & 96 & 128 & 192 \\
Missing endpoints / unknown labels & 0 / 0 & 0 / 0 & 0 / 0 \\
Applicable / inapplicable labels & 1,152 / 5,248 & 1,024 / 5,376 & 896 / 5,504 \\
Common fully labelled queries & 200 & 200 & 200 \\
Queries with at least two full endpoints & 200 & 200 & 200 \\
Queries with comparable effect pairs & 196 & 192 & 186 \\
Queries with an applicable candidate & 100 & 80 & 70 \\
Applicable-pool coverage (\%) & 50.0 & 40.0 & 35.0 \\
\bottomrule
\end{tabular}
\end{table}

\paragraph{Candidate-specific effect error.}
For the shared complete-endpoint candidate set $\mathcal V_t$ at query $t$, center each reference $F_{t,i}\in\{E_i,r_t,\widehat E_i\}$ and its observed target within that same set:
\begin{equation}
 \widetilde F_{t,i}=F_{t,i}-\frac{1}{|\mathcal V_t|}\sum_{j\in\mathcal V_t}F_{t,j},\quad
 \widetilde E^{\mathrm{obs}}_{t,i}=E^{\mathrm{obs}}_{t,i}
 -\frac{1}{|\mathcal V_t|}\sum_{j\in\mathcal V_t}E^{\mathrm{obs}}_{t,j}.
 \label{eq:rebuttal_centering}
\end{equation}
The centered error is $|\mathcal V_t|^{-1}\sum_{i\in\mathcal V_t}d_E(\widetilde F_{t,i},\widetilde E^{\mathrm{obs}}_{t,i})$, with $d_E$ from Eq.~\eqref{eq:distances}; queries with fewer than two complete endpoints are excluded from this metric and counted separately.
Because the centered query-only prediction is zero, this diagnostic tests candidate differences rather than the common query component.
Ordinary MSE and ordering agreement retain the definitions in Appendix~\ref{app:effect_probes}.
Development-branch repeatability checks fix ordering tolerances before held-out evaluation.

\begin{table}[H]
\centering\small
\caption{\textbf{Effect prediction within shared queries.}
Centering removes each query's mean effect. Macro averages weight tasks equally; query-only ordering is exactly $1/2$ on comparable pairs by definition, not an empirical gain.}
\label{tab:rebuttal_effects}
\begin{tabular}{@{}llccc@{}}
\toprule
Task & Effect reference & MSE $\downarrow$ & Centered MSE $\downarrow$ & Ordering $\uparrow$ \\
\midrule
PBB & Historical $E_i$ & 0.050 & 0.036 & 0.430 \\
 & Query-only $r_t$ & 0.039 & 0.035 & 0.500 \\
 & Predicted $\widehat E_i$ & 0.028 & 0.021 & 0.670 \\
PB & Historical $E_i$ & 0.058 & 0.039 & 0.410 \\
 & Query-only $r_t$ & 0.043 & 0.039 & 0.500 \\
 & Predicted $\widehat E_i$ & 0.033 & 0.024 & 0.640 \\
BT & Historical $E_i$ & 0.063 & 0.044 & 0.380 \\
 & Query-only $r_t$ & 0.048 & 0.043 & 0.500 \\
 & Predicted $\widehat E_i$ & 0.037 & 0.029 & 0.600 \\
\midrule
Macro & Historical $E_i$ & 0.057 & 0.040 & 0.407 \\
 & Query-only $r_t$ & 0.043 & 0.039 & 0.500 \\
 & Predicted $\widehat E_i$ & 0.033 & 0.025 & 0.637 \\
\bottomrule
\end{tabular}
\end{table}

\paragraph{Separate retrieval coverage from selection quality.}
For selection, we use one common set $\mathcal Q$ of queries whose entire valid candidate pool has known applicability labels, including known early-terminal outcomes.
Candidates are not removed because their predicted scores are poor or their $H$-step endpoints are unavailable.
Queries with unknown labels remain in the accounting; we additionally report all-query lower/upper bounds by assigning unknown outcomes 0/1.
Let $p^s_{t,i}$ be the selection probability for rule $s$, uniform over its maximum-score ties (absolute score tolerance $10^{-8}$), and let $\mathcal Q^+=\{t\in\mathcal Q:\max_i a_{t,i}=1\}$.
Define per-query choice accuracy $c_t^s=\sum_i p^s_{t,i}a_{t,i}$, pool coverage $\Pr(t\in\mathcal Q^+)$, all-query accuracy $\mathbb E[c_t^s]$, and covered-pool accuracy $\mathbb E[c_t^s\mid t\in\mathcal Q^+]$.
We use equal episode weight within each task, assigning equal weight to queries within episode; the conditional metric renormalizes those same query weights on $\mathcal Q^+$.
Within each task, all-query accuracy equals coverage times conditional accuracy, separating retrieval coverage from candidate selection quality.

The five rules are uniform random selection; historical compatibility $\kappa(E_i,r_t)$; predicted compatibility $\kappa(\widehat E_i,r_t)$; contextual selection $\ell_i^{\mathrm{rank}}+b_{t,i}$; and the full selector in Eq.~\eqref{eq:selection}.
They share the actual candidate pool, adapted actions, $r_t$, rank scores, and continuity state; only the scoring rule changes, and no counterfactual selection updates the cached cursor.
Removing $\kappa_i$ here is a \emph{fixed-checkpoint selector diagnostic}, not the separately configured no-comparison policy in Table~\ref{tab:ablation}.

\begin{table}[H]
\centering\footnotesize
\setlength{\tabcolsep}{4pt}
\caption{\textbf{Independent applicability of selected actions (\%).}
Each cell is all-query / covered-pool accuracy. All rules use the same labelled pools; the macro conditional column averages the task-conditional rates, not a pooled ratio.}
\label{tab:rebuttal_selection}
\begin{tabular}{@{}lcccc@{}}
\toprule
Selection rule & PBB & PB & BT & Macro \\
\midrule
Uniform random & 18.0 / 36.0 & 16.0 / 40.0 & 14.0 / 40.0 & 16.0 / 38.7 \\
Historical compatibility & 20.0 / 40.0 & 17.0 / 42.5 & 15.0 / 42.9 & 17.3 / 41.8 \\
Predicted compatibility & 23.0 / 46.0 & 20.0 / 50.0 & 17.0 / 48.6 & 20.0 / 48.2 \\
$\ell^{\mathrm{rank}}+b$ & 27.0 / 54.0 & 23.0 / 57.5 & 19.0 / 54.3 & 23.0 / 55.3 \\
Full selector & 31.0 / 62.0 & 28.0 / 70.0 & 23.0 / 65.7 & 27.3 / 65.9 \\
\bottomrule
\end{tabular}
\end{table}

\paragraph{Paired uncertainty and interpretation.}
We resample source episodes within each task 10,000 times, keeping all queries, branches, and scoring rules of each sampled episode together.
We report percentile 95\% intervals for paired differences and the equal-task mean, treating source episodes rather than candidate pairs as independent units.
Table~\ref{tab:rebuttal_effect_deltas} reports the principal contrasts, with task-wise directions visible in the preceding tables.

\begin{table}[H]
\centering\small
\caption{\textbf{Paired mechanism contrasts.}
Brackets give episode-clustered bootstrap 95\% intervals.
Centered-MSE contrasts are predicted minus reference; choice contrasts are percentage points.}
\label{tab:rebuttal_effect_deltas}
\begin{tabularx}{\linewidth}{@{}Xr@{}}
\toprule
Contrast & Estimate [95\% interval] \\
\midrule
Centered MSE: predicted minus query-only & $-0.014\ [-0.019,-0.009]$ \\
Centered MSE: predicted minus historical & $-0.015\ [-0.020,-0.009]$ \\
All-query choice: predicted minus historical compatibility & $+2.7\ [0.9,4.4]$ \\
All-query choice: full selector minus $\ell^{\mathrm{rank}}+b$ & $+4.3\ [2.0,6.7]$ \\
\bottomrule
\end{tabularx}
\end{table}

\paragraph{Candidate discrimination and decision quality.}
The predicted-versus-query-only centered-error difference is -0.014, while the full selector changes all-query choice accuracy by +4.3 percentage points relative to $\ell^{\mathrm{rank}}+b$ (Table~\ref{tab:rebuttal_effect_deltas}).
These contrasts distinguish two claims: sensitivity to candidate-specific consequences, and a benefit of that sensitivity for choosing actions.
A lower centered error together with a positive paired choice contrast supports both, whereas lower effect error alone did not establish better selection.
All-query accuracy (27.3\%) is interpreted jointly with pool coverage (41.7\%) and covered-pool accuracy (65.9\%), separating unavailable experience from selection errors.
These are pre-refinement diagnostics, not rollout success rates or evidence that every rejected proposal is useless after refinement.

\subsubsection{Adaptive Utility Beyond Deterministic Admission Rules}
\label{app:rebuttal_gate}

This study complements Table~\ref{tab:ablation} and Figure~\ref{fig:experience_distraction} by asking whether learned utility regulates recalled experience in practice, beyond the structural zero-gate property tested in Appendix~\ref{app:gate_probes}.

\paragraph{Record quality and gate values before refinement.}
We evaluate normal retrieval and the same random/similar complete-event replacements as Appendix~\ref{app:distraction_protocol} on PBB and PB.
For each condition, we prespecify ten replans from each of 20 held-out episodes per task and log the full learned-gate policy's selected proposal, raw $\alpha_t$, effective $g_t$, and all rule flags before executing the refined action.
We restore each sampled state to execute the selected adapted proposal and label its applicability using the independent task predicates above.
No training-negative/null flag is artificially enabled during evaluation.
Empty pools have $g_t=0$ but no selected-candidate applicability label and are recorded separately, without assigning an arbitrary $\alpha_t$ or negative label.
Unknown branch labels likewise remain in coverage counts.
Empty-pool / unknown-label counts are 0/0 for normal, 0 for random, and 0 for similar retrieval.

To separate learning from rules, define $v_t^{\mathrm{det}}=1$ when the selected proposal passes availability/provenance, temporal eligibility, and the hard repetition cutoff, \emph{before} the learned threshold $\alpha_{\min}$.
Raw-utility summaries and learned rejection rates are conditioned on $v_t^{\mathrm{det}}=1$; effective-gate summaries include deterministic vetoes for every labelled selected proposal.
This prevents a lower mean gate caused by hard rules from being attributed to learned discrimination.

\begin{table}[H]
\centering\footnotesize
\setlength{\tabcolsep}{3pt}
\caption{\textbf{Gate response to selected-proposal quality.}
Medians [Q1, Q3] replace a distribution plot. $N/N_{\mathrm{det}}$ gives labelled/eligible counts.
The final column is $\Pr(\alpha_t<\alpha_{\min}\mid v_t^{\mathrm{det}}=1,a)$ in percent, with the deployed threshold unchanged.
Empty pools and unknown labels are reported separately.}
\label{tab:rebuttal_gate_quality}
\begin{tabular}{@{}llrccc@{}}
\toprule
Retrieval & Applicability & $N/N_{\mathrm{det}}$ & $\alpha_t\mid v^{\mathrm{det}}=1$ & $g_t$ & Rejected (\%) \\
\midrule
Normal & Applicable & 260 / 252 & .72 [.56, .84] & .55 [.38, .68] & 4.0 \\
 & Inapplicable & 140 / 116 & .23 [.10, .39] & .04 [.00, .17] & 34.5 \\
Random & Applicable & 220 / 208 & .70 [.54, .82] & .53 [.35, .66] & 4.8 \\
 & Inapplicable & 180 / 140 & .19 [.08, .33] & .00 [.00, .13] & 46.4 \\
Similar & Applicable & 200 / 188 & .68 [.52, .81] & .51 [.32, .65] & 5.3 \\
 & Inapplicable & 200 / 148 & .14 [.05, .28] & .00 [.00, .08] & 54.1 \\
\bottomrule
\end{tabular}
\end{table}

\begin{table}[H]
\centering\small
\setlength{\tabcolsep}{4pt}
\caption{\textbf{Why candidate influence becomes zero (\%).}
First-veto attribution uses the displayed left-to-right order, so causes are mutually exclusive.
All rates use the full logged-query cohort per condition; the complete overlapping flag vector is retained in the logs.
Repetition's continuous attenuation is reflected in $g_t$, not counted as a hard veto.}
\label{tab:rebuttal_veto}
\begin{tabular}{@{}lccccc@{}}
\toprule
Retrieval & Other eligibility & Continuity & Repetition & $\alpha<\alpha_{\min}$ & Active $g>0$ \\
\midrule
Normal & 1.0 & 4.0 & 3.0 & 12.5 & 79.5 \\
Random & 2.0 & 6.0 & 5.0 & 18.8 & 68.2 \\
Similar & 3.0 & 7.0 & 6.0 & 22.5 & 61.5 \\
\bottomrule
\end{tabular}
\end{table}

The tables condition on the \emph{selected} proposal: a distractor-rich pool does not require a lower gate if selection still finds useful experience.
Moreover, $\alpha_t$ is trained as imitation utility, not a calibrated probability that a raw proposal is applicable.

\paragraph{A matched constant-utility control.}
We train a control that substitutes a development-selected constant $\alpha_0$ for the learned utility in \emph{both} source and conditioning, while retaining identical deterministic eligibility, continuity, hard/soft repetition, retrieval, selector architecture, losses, and training budget.
Candidate availability and eligibility checks, including the comparison to $\alpha_{\min}$, operate on the substituted value.
We choose $\alpha_0$ from $\{0.25,0.50,0.75\}$ using short, equally budgeted development-only runs and the equal-task/equal-condition mean SR; ties choose the smaller value.
This choice is frozen before full-budget training and held-out evaluation; the selected constant is $\alpha_0=\text{0.50}$.
The control retains the utility head and its auxiliary supervision but disconnects its output from action generation, isolating use of adaptive utility rather than removal of a loss or module.
Initialization, data ordering, optimizer schedules, NFE, action horizon/prefix, corruption schedules, and checkpoint-selection rules are matched.
Each variant is trained under its own deployed gate rule with three paired training seeds (11, 22, 33).
The retrieval/selector procedure is unchanged, but its fitted weights and visited states may differ across trained policies; only the fixed-query diagnostic above holds numerical scores and candidates identical.

Each trained checkpoint is evaluated on 100 paired resets per task/condition, with identical donor-sampling and policy-noise schedules at reset and the existing task step limits.
This gives 3,600 evaluation episodes across two variants, three seeds, two tasks, and three conditions, excluding development selection.
Subsequent trajectories are allowed to diverge; no shared on-policy candidate sequence is assumed.
We report mean $\pm$ sample SD across \emph{training seeds}, after averaging reset outcomes within seed.
This control differs from the no-comparison model, and its training-seed variability differs from the five evaluation-run variability in Figure~\ref{fig:experience_distraction}.

\begin{table}[H]
\centering\small
\setlength{\tabcolsep}{4pt}
\caption{\textbf{Learned versus constant utility (\%).}
Mean $\pm$ sample SD uses three training seeds; macro scores average both tasks within each seed before computing SD.
Drop is normal minus similar-distraction SR, computed within seed.}
\label{tab:rebuttal_constant_gate}
\begin{tabular}{@{}llcccc@{}}
\toprule
Task & Utility & Normal & Random & Similar & Similar drop $\downarrow$ \\
\midrule
PBB & Learned & $99.3\pm0.6$ & $94.0\pm2.0$ & $89.0\pm1.0$ & $10.3\pm0.6$ \\
 & Constant & $99.0\pm1.0$ & $89.7\pm1.5$ & $78.7\pm2.1$ & $20.3\pm1.5$ \\
PB & Learned & $92.0\pm1.0$ & $86.3\pm1.5$ & $81.3\pm1.2$ & $10.7\pm1.5$ \\
 & Constant & $91.3\pm1.5$ & $80.7\pm2.1$ & $69.7\pm2.1$ & $21.7\pm1.2$ \\
\midrule
Macro & Learned & $95.7\pm0.8$ & $90.2\pm1.4$ & $85.2\pm1.0$ & $10.5\pm1.0$ \\
 & Constant & $95.2\pm1.0$ & $85.2\pm0.6$ & $74.2\pm0.6$ & $21.0\pm0.5$ \\
\bottomrule
\end{tabular}
\end{table}

For uncertainty, we retain paired reset outcomes within task and seed and report the three seed-level contrasts separately.
A descriptive $t$ interval on paired seed differences uses $\bar d\pm t_{0.975,2}s_d/\sqrt3$; with only three seeds, this summary depends on a normality assumption and has limited precision.
The similar-condition macro contrast is 11.0 points (three seed-level contrasts: 11.5,9.5,12.0; descriptive 95\% interval: [7.7,14.3]).
The reduction in the normal-to-similar drop is 10.5 points (95\% interval: [8.3,12.7]).

\paragraph{Adaptive strength versus structural withdrawal.}
The matched control distinguishes the value of adaptive strength from the shared rules that disable candidate influence.
Comparable normal-condition success together with a smaller distraction-induced drop supports a benefit beyond those unchanged rules.
Higher raw utility for applicable proposals among deterministically eligible cases additionally localize that response to learning rather than veto frequency.
The closed-loop result is essential: low raw-proposal applicability alone does not prove that refinement cannot make use of the recalled structure.
The structural withdrawal guarantee alone does not establish a benefit of learned adaptive strength; that claim depends on the matched closed-loop comparison.
